\section{The Field Beneath the World}

We now follow the model into the parts of reality people have long called inner, hidden, or spiritual: contact between minds, messengers, stores of knowledge, other beings, and the largest question of all. The usual instinct is to file these as a soft epilogue after the hard physics. Resist that instinct, because in this model it is backwards. The foundation is experience, the vibes. The physical world is not the bedrock. It is one face of the field, the crystallized, shared face, downstream of the vibes and no more fundamental than the inner or the unseen faces. So these chapters are not a lesser appendix to the physics. They are the same one model, followed into its other faces, standing on the same ground.

\subsection{On status, honestly and evenly}

There is one real difference to be straight about, and it is not a difference of legitimacy. The simulation built so far has mostly been pointed at the physical results, so the inner and field aspects of the model have not been run on the bench in the same direct way. But that is a matter of where the work has gone, not of what is real, and plenty of the physics is also unfinished. The physics holds no special authority here. It is itself only patterns in the field, a derived face, with no standing to rule on what is real at the foundation, which is experience. So the chapters that follow draw from the one model exactly as the physics does, built from the same pieces, and they are presented as real parts of the picture. Where something has not been directly demonstrated, that is said plainly, the same as anywhere else in the walkthrough, no more and no less.

\subsection{The ether is the field below the shared pattern}

Recall the two conditions, lawful and wild. The physical world is the lawful, shared region of the field, the part that has crystallized into one stable pattern everyone inhabits. Beyond it lies the rest of the field, the part not bound up in that shared physical pattern. Some of that rest is wild and fluid, like raw feeling. Some of it is lawful in its own right, structured and stable but not physical. What older language calls the ether is, in these terms, simply the field itself beneath the shared physical appearance. It is not a separate realm bolted on. It is the same one web, viewed below the level of the physical pattern.

\subsection{The mechanism, built from pieces already in hand}

Everything that follows uses only parts the model already contains, and it is worth setting out the pieces once so the later chapters can lean on them.

There is one shared substrate, so two selves are never wholly separate. This is the very same fact behind quantum entanglement. Influence travels by notes, paths through the field that connect two structures. A path carries a pattern well when its two ends are in tune, resonant, the way a radio passes a station only when matched to its frequency. Once a path is open and resonant, the one rule carries the pattern along it. And the arriving pattern is unpacked by the receiver's own sub-selves, becoming words for some, a wordless impression for others.

One further fact does the heavy lifting, and it is what makes any of this conceivable: distance in the field is not the same as distance in physical space. The substrate is tree-like, and in a tree two points far apart along the surface can be close through a shared deep root. So two selves far apart in physical space can be near in the field if they share deep structure. Contact runs through that shared depth, not through the air between two bodies.

\subsection{How fast, and how far}

One thing all of it shares is how it travels. It travels by the one rule, at the rate of the beat, stepping along paths in the field. That is not instantaneous magic. But the rate of the beat is not obviously the same as the speed of light. Light is the speed at which influence crosses the physical surface of the crystal, and these paths run instead through the field's depth, where distance is not physical distance. Two points far apart across the physical surface can be only a few steps apart through a shared deep root, so a beat-paced signal along that deep path could arrive far sooner than light crossing the physical space between them. Whether such contact actually outruns physical light we do not claim to know, and we drop the assumption that the speed of light must bound it. What looks like contact across a distance would be a shared depth surfacing on two ends, traveling the short way through the field rather than the long way through space.

With that foundation set, the next chapters take the cases in turn: mind-to-mind contact, messengers and channeling, hidden stores of knowledge, other realms, and the question of God. Each is a particular use of this one mechanism, and each belongs to the model as fully as the physics does.

\textbf{Takeaway:} beneath the shared physical world lies the rest of the field, the ether, and anything in it can reach anything else by one mechanism: a shared substrate, paths along notes, resonance to open a path, propagation by the rule, and the fact that field-distance is not physical distance, so depth makes the distant near. It travels by the beat along field-paths that can be far shorter than the physical route, so whether it outruns physical light is left open. None of this is a lesser tier than the physics, since experience, not physics, is the foundation.
